\section{Mục đích}
Tài liệu này mô tả các yêu cầu phần mềm của hệ thống NaviMart. Nội dung tài liệu tập trung vào mục tiêu xây dựng hệ thống, phạm vi chức năng, các tác nhân sử dụng, thuật ngữ nghiệp vụ và các tài liệu tham khảo liên quan.

Tài liệu được dùng làm cơ sở thống nhất yêu cầu giữa các bên liên quan, nhóm phát triển, nhóm kiểm thử và người vận hành hệ thống trong suốt quá trình phân tích, thiết kế, xây dựng, kiểm thử và bảo trì phần mềm.

\section{Phạm vi}
NaviMart là hệ thống hỗ trợ người dùng quản lý việc mua sắm và sử dụng thực phẩm trong gia đình. Hệ thống cho phép người dùng tạo danh sách mua sắm, quản lý thực phẩm trong tủ lạnh/kho thực phẩm, lập kế hoạch bữa ăn, tìm kiếm/gợi ý công thức nấu ăn, theo dõi thống kê tiêu dùng và lãng phí thực phẩm.

Mục tiêu chính của NaviMart là tạo ra một luồng nghiệp vụ khép kín: mặt hàng trong danh sách mua sắm sau khi được đánh dấu đã mua sẽ được chuyển vào kho thực phẩm; dữ liệu kho thực phẩm được sử dụng để gợi ý món ăn, lập kế hoạch bữa ăn và kiểm tra nguyên liệu; các nguyên liệu còn thiếu có thể được thêm ngược lại vào danh sách mua sắm.

Hệ thống được tổ chức thành các module chính sau:
\begin{itemize}
    \item \textbf{Xác thực và quản lý tài khoản:} Hỗ trợ đăng ký, đăng nhập, đăng nhập qua mạng xã hội, quên mật khẩu, đổi mật khẩu, đăng xuất và cập nhật thông tin cá nhân.
    \item \textbf{Quản lý nhóm gia đình:} Hỗ trợ tạo nhóm gia đình, tham gia nhóm, rời khỏi nhóm, xem thông tin nhóm, tạo mã mời và cập nhật phân quyền.
    \item \textbf{Quản lý danh sách mua sắm:} Hỗ trợ tạo danh sách, thêm/cập nhật/xóa mặt hàng và đánh dấu mặt hàng đã mua.
    \item \textbf{Quản lý tủ lạnh và kho thực phẩm:} Hỗ trợ xem, tìm kiếm, thêm, cập nhật, thanh lý thực phẩm, tự động trừ kho khi nấu ăn và thông báo thực phẩm sắp hết hạn.
    \item \textbf{Lên kế hoạch bữa ăn:} Hỗ trợ xem lịch trình, tạo phiên ăn, thêm/sửa/xóa món ăn trong kế hoạch.
    \item \textbf{Gợi ý và công thức nấu ăn:} Hỗ trợ tìm kiếm công thức, xem chi tiết công thức, lưu/bỏ lưu công thức, gợi ý món ăn theo Pantry, chat với AI Chef và thêm nguyên liệu thiếu vào danh sách mua sắm.
    \item \textbf{Báo cáo và thống kê:} Hỗ trợ người dùng xem thống kê tiêu dùng và lãng phí thực phẩm.
    \item \textbf{Quản trị hệ thống:} Hỗ trợ Quản trị viên xem danh sách user, khóa/mở khóa tài khoản, quản lý master data, kiểm duyệt công thức và xem dashboard tổng quan.
\end{itemize}

Các tác nhân của hệ thống gồm ba tác nhân chính và hai tác nhân hỗ trợ tự động:
\begin{itemize}
    \item \textbf{Khách:} Người chưa có tài khoản hoặc chưa đăng nhập, có thể đăng ký tài khoản và đăng nhập qua mạng xã hội.
    \item \textbf{Người dùng:} Người có tài khoản trong hệ thống, sử dụng các chức năng phục vụ quản lý mua sắm, kho thực phẩm, bữa ăn, công thức và thống kê.
    \item \textbf{Quản trị viên:} Người vận hành hệ thống, quản lý người dùng, dữ liệu nền, nội dung công thức và dashboard tổng quan.
    \item \textbf{Hệ thống:} Tác nhân hỗ trợ tự động, thực hiện các xử lý nền như chuyển mặt hàng đã mua vào kho, trừ kho khi nấu ăn và thông báo thực phẩm sắp hết hạn.
    \item \textbf{Hệ thống AI:} Tác nhân hỗ trợ tự động, tham gia vào chức năng gợi ý món ăn theo Pantry và phản hồi trong chức năng AI Chef.
\end{itemize}

\section{Từ điển thuật ngữ}
\begin{itemize}
    \item \textbf{NaviMart:} Tên hệ thống hỗ trợ quản lý mua sắm, kho thực phẩm, kế hoạch bữa ăn, công thức và thống kê tiêu dùng thực phẩm.
    \item \textbf{Khách:} Người chưa có tài khoản hoặc chưa đăng nhập vào hệ thống.
    \item \textbf{Người dùng:} Người có tài khoản trong NaviMart và có thể sử dụng các chức năng chính của hệ thống.
    \item \textbf{Quản trị viên:} Người có quyền quản trị, vận hành và kiểm soát các dữ liệu dùng chung của hệ thống.
    \item \textbf{Nhóm gia đình:} Không gian dùng chung để các tài khoản người dùng phối hợp quản lý danh sách mua sắm, kho thực phẩm và kế hoạch bữa ăn.
    \item \textbf{Quyền trong nhóm gia đình:} Cấp quyền nội bộ của Người dùng trong một nhóm gia đình, ví dụ Owner, Editor hoặc Viewer. Đây là quyền sử dụng chức năng trong nhóm, không phải tác nhân độc lập của hệ thống.
    \item \textbf{Danh sách mua sắm:} Danh sách các mặt hàng cần mua. Mặt hàng có thể được đánh dấu là đã mua và chuyển vào kho thực phẩm.
    \item \textbf{Tủ lạnh/Kho thực phẩm/Pantry:} Nơi hệ thống lưu trữ dữ liệu thực phẩm hiện có, bao gồm tên, số lượng, hạn sử dụng và thông tin liên quan.
    \item \textbf{Kế hoạch bữa ăn:} Lịch trình các phiên ăn và món ăn dự kiến theo ngày hoặc theo tuần.
    \item \textbf{Công thức nấu ăn:} Thông tin hướng dẫn nấu ăn, bao gồm nguyên liệu và các bước thực hiện.
    \item \textbf{AI Chef:} Chức năng chat/gợi ý thông minh, hỗ trợ người dùng lựa chọn món ăn dựa trên nhu cầu và dữ liệu Pantry.
    \item \textbf{Hệ thống AI:} Tác nhân hỗ trợ tự động đại diện cho thành phần AI trong các use case gợi ý món ăn và chat AI Chef.
    \item \textbf{Master Data:} Dữ liệu nền dùng chung như danh mục thực phẩm, đơn vị tính và các thông tin chuẩn hóa khác.
    \item \textbf{Dashboard:} Giao diện tổng hợp các chỉ số vận hành và thống kê tổng quan dành cho Quản trị viên.
\end{itemize}

\section{Tài liệu tham khảo}
\begin{itemize}
    \item Lecture note.
    \item Biểu đồ use case và activity diagram của hệ thống NaviMart.
\end{itemize}
